\begin{abstract}
Motion Manifold Primitives (MMP), a manifold-based approach for encoding basic motion skills, can produce diverse trajectories, enabling the system to adapt to unseen constraints. Nonetheless, we argue that current MMP models lack crucial functionalities of movement primitives, such as temporal and via-points modulation, found in traditional approaches. This shortfall primarily stems from MMP's reliance on discrete-time trajectories.
To overcome these limitations, we introduce Motion Manifold Primitives++ (MMP++), a new model that integrates the strengths of both MMP and traditional methods by incorporating parametric curve representations into the MMP framework. Furthermore, we identify a significant challenge with MMP++: performance degradation due to geometric distortions in the latent space, meaning that similar motions are not closely positioned. To address this, Isometric Motion Manifold Primitives++ (IMMP++) is proposed to ensure the latent space accurately preserves the manifold's geometry.
Our experimental results across various applications, including 2-DoF planar motions, 7-DoF robot arm motions, and SE(3) trajectory planning, show that MMP++ and IMMP++ outperform existing methods in trajectory generation tasks, achieving substantial improvements in some cases. Moreover, they enable the modulation of latent coordinates and via-points, thereby allowing efficient online adaptation to dynamic environments.

% A manifold representation of basic motion skills, referred to as Motion Manifold Primitives (MMP), has recently been investigated as highly adaptable models capable of generating diverse trajectories that can complete the given task. However, in this paper, we claim that existing MMP methods lack important functionalities of movement primitives found in other conventional methods (e.g., temporal modulation, via-points modulation, etc.). This limitation originates from their reliance on discrete-time trajectory representations, while others employ parametric curve models with special inductive biases.
% We propose {\it Motion Manifold Primitives++ (MMP++)}, which combines the advantages of the MMP and conventional methods by applying the MMP framework to the parametric curve representations. 
% In addition, we observe that the performance of MMP++ can sometimes degrade significantly due to geometric distortion in the latent space -- by distortion, we mean that similar motions are not located nearby in the latent space. To mitigate this issue, we propose {\it Isometric Motion Manifold Primitives++ (IMMP++)}, where the latent coordinate space is regularized to preserve the geometry of the manifold.
% Experimental results with 2-DoF planar motions, 7-DoF robot arm movements, and SE(3) trajectory generation tasks demonstrate that MMP++ and IMMP++ enable intuitive modulation of latent coordinates and via-points, leading to online adaptation in dynamic environments.
\end{abstract} 